\chapter{La página web, en paralelo a todas las fases}
\label{cap:d-pagina}
\textbf{Fechas}: del 28 de septiembre al 2 de octubre de 2026. \textbf{Notas}: \codigo{06-pagina.md}, \codigo{07-diseno.md} y
de la 12 a la 17.

La página se construyó junto con los datos (decisión 2.1). Estas son las decisiones que la fueron formando, en orden.

\begin{opciones}[P.1 · Para gente no técnica · 28-sep-2026]
\textbf{Lo que pidió Brandon}: enfocarse en la página, ``para gente no técnica''.

\textbf{Decisiones}: una sola idea por sección, cifras medidas y en palabras (``de cada 10 cuartos, 6 ocupados'').
\textbf{Descartado}: destacar en la portada a Kohunlich (+13.5\,\%) y Dzibanché (+69.2\,\%), porque la Ruta completa bajó
8.2\,\% (escoger datos a conveniencia), y una gráfica de ``ocupación hotelera'' de cuatro centros que pedía conocimientos
técnicos.

\textbf{Ubicación comprobada}: los 11 pueblos están en el Censo de Quintana Roo, las 4 zonas tienen ``Quintana Roo'' en el INAH
y cada coordenada cae en el polígono de su municipio. No basta el nombre: hay un ``Xul-Ha'' en Puerto Morelos y un ``Felipe
Carrillo Puerto'' en Solidaridad.
\end{opciones}

\begin{opciones}[P.2 · El sistema visual ``Sur mexicano'' · 28-sep-2026]
\textbf{Lo que dijo Brandon}: \emph{``se ve muy IA el font y todo''}, \emph{``veo mucho ruido e información innecesaria, no es nada
visual''}, \emph{``calidad tipo Apple, hermoso''}, \emph{``tipografía de destino, viajes, mar, mexicano\dots{} colores vibrantes para
México, identidad real''}.

\textbf{Elegido}: rosa mexicano, amarillo cempasúchil, turquesa y añil en bloques de color completo; letras Bricolage
Grotesque y Figtree guardadas en la propia página (funcionan sin internet); la greca escalonada con el perfil de las pirámides
del sur; fotos grandes y unos 70\,\% menos de texto.

\textbf{Descartados}: el estilo del primer diseño (genérico, ``de IA''); imitar literalmente a Apple (se tomaron sus
principios, no su diseño); el papel picado (cliché nacional; la greca es propia del sur de Quintana Roo).
\end{opciones}

\begin{opciones}[P.3 · La comida, y por qué se retiró · 1-oct-2026]
\textbf{Lo que pidió Brandon}: \emph{``Agrega muchisimas fotos de todo la comida de la vista nace el amor sabes''}.

\textbf{Lo que se hizo}: una galería de 24 platillos de la península con fotos de Wikimedia Commons, con licencia libre y
crédito. Se revisaron a ojo y se rechazaron, por ejemplo, alfajores argentinos que salieron buscando ``relleno negro''.

\textbf{El error}: solo se había revisado que las fotos no mencionaran lugares excluidos. Eran de \textbf{Mérida, Campeche y
otros lugares fuera de Quintana Roo}.
\end{opciones}

\begin{opciones}[P.4 · Fotos comprobadas y una página que empieza por el viaje · 1-oct-2026]
\textbf{Lo que dijo Brandon}: \emph{``pero que sean fotos de las zonas que si sean de quintana rooo\dots{} tambien necesito foto
dependiendo que escojan. Y quiero primero hasta arriba en el inicio donde escojas como el plan de viaje''}.

\textbf{Decisiones}:
\begin{itemize}
\item La galería de platillos se \textbf{retiró completa}.
\item \textbf{Toda foto debe tener coordenada GPS dentro del municipio de su lugar}; si no, el programa se detiene. De 331 fotos
  con coordenada se eligieron 28, revisadas a ojo.
\item La foto de la portada cambia con el lugar elegido.
\item Nuevo orden: el planeador arriba, ``Así va a estar'', ``Qué hacer'', los lugares, y ``Los datos'' al final.
\end{itemize}
\textbf{Huecos declarados}: en Commons no hay fotos de platillos tomadas en los 5 lugares (0 de 331), y Rest-Mex no tiene
reseñas de ellos. Mostrarlas sería inventarlas.
\end{opciones}

\begin{opciones}[P.5 · Una página que convence · 2-oct-2026]
\textbf{Lo que pidió Brandon}: \emph{``más fotos a las comidas, reseñas reales, las estrellas del lugar\dots{} preguntas frecuentes,
secciones para vender más cosas, una barra horizontal de imágenes bonitas\dots{} en 10 idiomas, una versión nocturna y otra de
día\dots{} sin saturar''}.

\begin{center}
\small
\begin{tabularx}{\linewidth}{>{\raggedright\arraybackslash}p{0.22\linewidth}YY}
\encabezado \h{Pedido} & \h{Lo que existe} & \h{Lo que eligió Brandon} \\
Reseñas reales & Rest-Mex no cubre los 5 lugares; copiar las de Google está prohibido & Botón ``Reseñas en Google'' en cada
negocio + reseñas del equipo \\
\filagris Estrellas & DataTur las publica solo para 70 centros, ninguno del sur & Estrellas oficiales donde existen; hueco
declarado en el sur \\
Secciones para vender & No hay precios oficiales & Experiencias reales, rutas de 2 y 3 días y llamados, sin precios \\
\filagris 10 idiomas & — & 9 publicados; el maya yucateco espera la revisión de un hablante; ``Los datos'' en español \\
Día y noche & — & Botón animado de sol a luna; la página se abre como un círculo \\
\bottomrule
\end{tabularx}
\normalsize
\end{center}
\end{opciones}

\begin{opciones}[P.6 · Cancún y Riviera Maya en toda la parte del viajero · 2-oct-2026]
\textbf{Lo que dijo Brandon}: \emph{``se te olvidó los cambios de meter cancún y riviera maya, como opciones, porque es que sigues
mencionando otros cuando no tenemos datos de ninguno de ellos, sabes? y luego el modo nocturno no sirve''}.

\textbf{Decisiones}: en el mapa, las fichas, las experiencias, las rutas y las preguntas, los lugares del viajero son Chetumal,
Calderitas–Oxtankah, la Ruta, Cancún y la Riviera Maya (estos dos siempre con su etiqueta de referencia). ``Los datos'' no
cambia: sigue siendo el análisis de las 5 regiones. El modo noche se rehízo: cada bloque de color tiene su versión oscura.
Las estrellas oficiales del norte (Cancún: 67.8\,\% de sus cuartos son de 5 estrellas) y 8 fotos de comida con coordenada
comprobada en Cancún y Playa del Carmen.
\end{opciones}
